%%=============================================================================
%% Conclusie en reflectie
%%=============================================================================

\chapter{\IfLanguageName{dutch}{Conclusie en reflectie}{Conclusion and reflection}}%
\label{ch:reflectie}
Het graduaatsproject richt zich op de ontwikkeling van een zelfstandige proof
of concept (PoC) van een visuele e-maileditor voor Vesalius.ai. Het project
combineert het analyseren van een technisch probleem met het onderzoeken,
ontwerpen en realiseren van een oplossing binnen een professionele
ontwikkelcontext.

De focus ligt op het visueel aanpassen van automatisch gegenereerde
e-mailtemplates, het veilig verwerken van de inhoud en het controleren van
de gegenereerde e-mails in gangbare e-mailclients. Daarbij wordt rekening
gehouden met de beperkingen van e-mailtechnologieën en met de noodzaak om
de oplossing gescheiden te houden van de productieomgeving.

\section{\IfLanguageName{dutch}{Conclusie}{Conclusion}}%
\label{sec:conclusie}
De doelstelling van het graduaatsproject is het ontwikkelen van een
zelfstandige PoC waarmee gebruikers de inhoud en ondersteunde opmaak van
bepaalde e-mailtemplates visueel kunnen aanpassen, opslaan, vooraf bekijken
en testen. Daarnaast wordt onderzocht hoe de gegenereerde HTML veilig kan
worden verwerkt en hoe de e-mails worden weergegeven in verschillende
e-mailclients.

De PoC richt zich op drie piloottypes: afspraakbevestigingen,
afspraakherinneringen en afspraakannuleringen. De vaste e-mailschil, met
bijvoorbeeld het logo, de merkkleuren en de footer, blijft grotendeels
behouden. De nadruk ligt op de bewerkbare inhoud en op een betrouwbare
verwerking van de uiteindelijke e-mail.

De meerwaarde voor Vesalius.ai bestaat erin dat het project de technische
haalbaarheid van visueel aanpasbare e-mailtemplates onderzoekt. Een
geslaagde PoC kan aantonen welke onderdelen nodig zijn om gebruikers meer
controle te geven over de inhoud van e-mails, zonder dat hiervoor onmiddellijk
een integratie in het bestaande productieplatform nodig is.

De uiteindelijke beoordeling van de doelstelling gebeurt aan de hand van
de gerealiseerde functionaliteiten, de geautomatiseerde tests en de
compatibiliteitscontroles. Op basis daarvan kan worden vastgesteld welke
vereisten volledig zijn gerealiseerd, welke slechts gedeeltelijk zijn
aangetoond en welke verdere ontwikkeling vereisen. De PoC vormt daarmee
een mogelijke technische basis voor een toekomstige integratie, maar is
zelf geen productieklare oplossing.

\section{\IfLanguageName{dutch}{Wat ik geleerd heb}{What I learned}}%
\label{sec:geleerd}
Een belangrijk leerpunt is het belang van een duidelijke afbakening bij
softwareontwikkeling. Een visuele e-maileditor lijkt op het eerste gezicht
een afgebakende gebruikersfunctionaliteit, maar omvat ook templatebeheer,
HTML-verwerking, dynamische gegevens, CSS en compatibiliteit met
verschillende e-mailclients. Door deze onderdelen afzonderlijk te
onderzoeken, kan de complexiteit beter worden ingeschat en kan de oplossing
stapsgewijs worden opgebouwd.

Daarnaast leer ik dat een technische keuze niet alleen op basis van
gebruiksgemak moet worden gemaakt. Een editor moet bijvoorbeeld niet enkel
voldoende bewerkingsmogelijkheden aanbieden, maar ook HTML genereren die
veilig kan worden verwerkt en geschikt is voor gebruik in e-mails. De
geschiktheid van een technologie moet daarom worden beoordeeld aan de hand
van concrete vereisten en technische tests.

Ook het onderscheid tussen een preview en de uiteindelijke weergave van
een e-mail is belangrijk. Een e-mail die correct wordt weergegeven in een
browser, wordt niet noodzakelijk identiek weergegeven in Gmail, Outlook of
Apple Mail. Dit maakt duidelijk waarom compatibiliteitscontroles een
afzonderlijk onderdeel van het ontwikkelproces moeten vormen.

Op professioneel vlak biedt het project de mogelijkheid om ervaring op te
doen met het opdelen van een probleem in afzonderlijke technische onderdelen,
het documenteren van ontwerpkeuzes en het systematisch controleren van
vereisten. De werkzaamheden worden opgevolgd via Jira en de broncode wordt
beheerd met Git en Bitbucket. De bestaande bedrijfscontext vraagt bovendien
aandacht voor vertrouwelijkheid en voor het correct omgaan met interne
broncode en configuratiegegevens.

De verdere reflectie wordt aangevuld op basis van de werkelijk uitgevoerde
werkzaamheden, de technische problemen die tijdens de ontwikkeling werden
vastgesteld en de resultaten van de uitgevoerde tests.

\section{\IfLanguageName{dutch}{Wat ik anders zou doen}{What I would do differently}}%
\label{sec:anders}
Bij een gelijkaardig project zou ik de technische risico's zo vroeg mogelijk
onderzoeken. Voor een WYSIWYG-e-maileditor betekent dit onder meer dat ik
vroegtijdig zou testen welke HTML de gekozen editor genereert, welke
bewerkingsmogelijkheden nodig zijn en welke beperkingen zich voordoen bij
de weergave in verschillende e-mailclients.

Daarnaast zou ik de functionele vereisten vanaf het begin vertalen naar
concrete testscenario's. Voorbeelden hiervan zijn het aanpassen en opnieuw
laden van een template, het veilig verwerken van placeholders en het
controleren van de gerenderde HTML. Hierdoor kan tijdens de ontwikkeling
sneller worden vastgesteld of een onderdeel aan de verwachtingen voldoet.

Ook zou ik ontwerpbeslissingen, technische risico's en openstaande vragen
systematisch documenteren. Dit maakt het eenvoudiger om achteraf te
verklaren waarom een bepaalde technologie of architectuur werd gekozen.
Het helpt bovendien om de voortgang te evalueren en om te vermijden dat
belangrijke beslissingen uitsluitend in gesprekken of persoonlijke notities
worden bijgehouden.

Ten slotte zou ik voldoende tijd reserveren voor testen in echte
e-mailclients. De implementatie van de editor en de renderingpipeline is
slechts een deel van het werk; de uiteindelijke bruikbaarheid hangt ook af
van de manier waarop de gegenereerde e-mails worden weergegeven.

\section{\IfLanguageName{dutch}{Verder werk}{Future work}}%
\label{sec:verder-werk}
Een mogelijke volgende stap is de integratie van de PoC in de bestaande
architectuur van Vesalius.ai. Hiervoor moet eerst worden onderzocht hoe de
oplossing aansluit op het bestaande templatebeheer, de e-mailverwerking,
de opslag en de geldende beveiligingsvereisten. Deze integratie valt buiten
de scope van het huidige graduaatsproject.

Daarnaast kan de compatibiliteit van de gegenereerde e-mails verder worden
onderzocht. Verdere tests in verschillende versies en platformen van Gmail,
Outlook en Apple Mail kunnen aantonen welke HTML- en CSS-elementen
betrouwbaar worden weergegeven en waar aanvullende aanpassingen nodig zijn.

Ook de editor zelf kan verder worden uitgebreid. Mogelijke uitbreidingen
zijn extra opmaakmogelijkheden, ondersteuning voor bijkomende e-mailtypes
en verbeteringen aan het beheer en de versieopvolging van templates.
Dergelijke uitbreidingen moeten worden afgewogen tegen de noodzaak om de
editor eenvoudig en beheersbaar te houden.

Ten slotte kan een toekomstig project onderzoeken hoe de PoC op een veilige
en onderhoudbare manier in de productieomgeving kan worden opgenomen.
Daarbij zijn onder meer de keuze van de definitieve opslagmethode, de
integratie met bestaande diensten, toegangsbeheer, validatie, monitoring
en het beheer van wijzigingen aan templates van belang.

De resultaten van dit graduaatsproject kunnen als uitgangspunt dienen om
de technische haalbaarheid, de resterende risico's en de vereisten voor
een mogelijke verdere ontwikkeling te beoordelen.